\documentclass[11pt]{article}
\usepackage{amsmath,amssymb,amsthm}
\newtheorem{theorem}{Theorem}
\newtheorem{lemma}{Lemma}
\title{Exact shared amplitudes at half load}
\author{}\date{}
\begin{document}\maketitle

\paragraph{Status.}
This is a complete written interpolation lemma, with an exact rational
checker. It is not Lean-checked. It supplies two requested root words
using one absolute residue pool and independently chosen signs. It
does not supply two arbitrary descendant tables and proves neither
accepted asymptotic HeadComplexity target.

Fix $q\geq1$, distinct real poles $p_1,\ldots,p_{2q}$, and distinct
real query points $t_1,\ldots,t_q$ different from every pole. Write
\[
 P=P_S P_T,\qquad
 P_S(t)=\prod_{h=1}^q(t-p_h),\qquad
 P_T(t)=\prod_{h=q+1}^{2q}(t-p_h).
\]

\begin{theorem}[A rectangular shared pool]
For every two real vectors $v,w\in\mathbb R^q$, there are polynomials
$N_1,N_2$ of degree less than $2q$ such that
\[
 \frac{N_1(t_i)}{P(t_i)}=v_i,\qquad
 \frac{N_2(t_i)}{P(t_i)}=w_i,
 \qquad N_2(p_h)=\begin{cases}-N_1(p_h)&h\leq q,\\
                              N_1(p_h)&h>q.\end{cases}       \tag{1}
\]
Consequently the two rational functions have shared absolute residues
at all $2q$ poles. For every two strict sign words on the query grid,
one can arrange for all these absolute residues to be strictly positive.
All data may be rational when the poles and queries are rational.
\end{theorem}
\begin{proof}
Interpolate polynomials $a,b$ of degree less than $q$ at the query
points, prescribing
\[
 a(t_i)=\frac{P_T(t_i)(v_i+w_i)}2,\qquad
 b(t_i)=\frac{P_S(t_i)(v_i-w_i)}2.
\]
Set $A=P_Sa$, $B=P_Tb$, $N_1=A+B$, and $N_2=A-B$.
Their degrees are less than $2q$, and direct substitution proves
(1). Since the poles are simple, partial fractions give
\[
 \frac{N_\ell(t)}{P(t)}
   =\sum_{h=1}^{2q}\frac{N_\ell(p_h)}{P'(p_h)(t-p_h)}.
\]
The last equality in (1) proves equality of the absolute residues.

For strict sign words, initially take $v_i,w_i\in\{-1,1\}$.
As $(v,w)$ varies in a sufficiently small open neighborhood, its
signs do not change. The interpolation construction is a linear
isomorphism from $(v,w)$ to the coefficient vectors of $(a,b)$.
For each $h\leq q$, $N_1(p_h)=P_T(p_h)b(p_h)$ is a nonzero
linear functional of these vectors. For each $h>q$, the same holds
for $N_1(p_h)=P_S(p_h)a(p_h)$. Avoiding these finitely many proper
hyperplanes makes every residue nonzero without changing any query
sign. Rational perturbations suffice for rational data, because a
finite union of proper rational hyperplanes cannot contain a rational
open box. Taking absolute values then gives a strictly positive pool.
\end{proof}

\paragraph{How this enters a coalescing root construction.}
For a signed incoming numerator $N$, the adaptive-pole calculation
from \texttt{adaptive-barycentric-root.tex} has the following exact
extension. Let $P(t)=\prod_s(t-x_s)$, choose $R$ of degree less
than the number of centers, and suppose $N(x_s)R(x_s)\ne0$.
Choose a mask sign $\eta_s$ so that
\[
 d_s=\frac{\eta_sN(x_s)}{R(x_s)P'(x_s)}>0.
\]
Put $p_s^\pm=x_s\pm d_s/\lambda$, $Q_\lambda=
\prod_s(t-p_s^-)(t-p_s^+)$, and
$c_s^\pm=N(p_s^\pm)/Q_\lambda'(p_s^\pm)$.
Then the masked score with coefficients
$(-\eta_s c_s^-,\eta_s c_s^+)$ satisfies
\[
 \frac1\lambda\sum_s\left(
 \frac{-\eta_s c_s^-}{t-p_s^-}
 +\frac{\eta_s c_s^+}{t-p_s^+}\right)
 \longrightarrow \frac{R(t)}{P(t)}.                         \tag{2}
\]
The child receiving the mask-positive pole inherits residues
$R(x_s)/(2P'(x_s))$ after division by $\lambda$; the other child
inherits their negatives. Indeed,
\[
 c_s^\pm/\lambda\longrightarrow
 \frac{\pm N(x_s)}{2d_sP'(x_s)^2},
\]
and substitution proves both assertions. Thus an incoming numerator
need not be constant. But the same $R$ still reaches both children;
changing the incoming numerator does not remove the compatibility
problem.

\paragraph{Scope and exact-value cost.}
The theorem escapes the square shared-pool obstruction by using twice
as many poles as queries. It permits the pool to depend on the two
requested words. It does not produce one pool for all words.
Its value interpolation uses exactly $2q$ continuous coefficients.
Applying it to all descendant queries would require counting all of
them against these coefficients. No logarithmic saving follows from
repeating this exact-value construction.

The checker verifies all ordered pairs of words for $1\leq q\leq4$,
using exact rational arithmetic. It perturbs the query values inside
their requested orthants, checks that every residue is nonzero, checks
the prescribed relative residue signs, and checks both partial
fraction scores at every query.
\end{document}
